\documentclass{article}
\usepackage{amsmath,amssymb,amsthm}
\usepackage{mathrsfs}
\usepackage{enumitem}

\newtheorem{theorem}{Theorem}[section]
\newtheorem{lemma}[theorem]{Lemma}
\newtheorem{proposition}[theorem]{Proposition}
\newtheorem{corollary}[theorem]{Corollary}
\theoremstyle{definition}
\newtheorem{definition}[theorem]{Definition}
\theoremstyle{remark}
\newtheorem{remark}[theorem]{Remark}
\newtheorem{example}[theorem]{Example}
\newtheorem{physremark}[theorem]{Physical remark}

% Notation
\newcommand{\ad}{\operatorname{ad}}
\newcommand{\End}{\operatorname{End}}
\newcommand{\tr}{\operatorname{tr}}
\newcommand{\K}{K}
\newcommand{\g}{\mathfrak{g}}
\newcommand{\h}{\mathfrak{h}}
\newcommand{\R}{\mathbb{R}}
\newcommand{\C}{\mathbb{C}}
\newcommand{\Z}{\mathbb{Z}}
\newcommand{\D}{\Delta}
\newcommand{\W}{\mathcal{W}}
\newcommand{\Span}{\operatorname{span}}
\newcommand{\rad}{\operatorname{rad}}

\title{A Comprehensive Introduction to Semisimple Lie Algebras\\
for Theoretical Physicists}
\author{}
\date{}

\begin{document}
\maketitle
\tableofcontents

\section*{Introduction}
Lie algebras are the natural language of continuous symmetries. In quantum mechanics,
the angular momentum operators $L_x, L_y, L_z$ close under commutation, forming the
Lie algebra $\mathfrak{so}(3)$. In particle physics, the Gell‑Mann matrices satisfy
the commutation relations of $\mathfrak{su}(3)$, the algebra of the strong force.
More generally, any continuous symmetry group of a physical system gives rise to a
Lie algebra of infinitesimal transformations, whose structure deeply constrains the
possible spectra, interactions, and dynamics.

This text aims to provide theoretical physicists with a self‑contained, rigorous
yet physically motivated exposition of semisimple Lie algebras and their
classification. We will follow the notation and logical development of Brian Hall’s
textbook \cite{Hall}, filling in physical insights, examples, and applications at
every stage. By the end, the reader will understand why the possible simple Lie
algebras are precisely $A_n, B_n, C_n, D_n, E_6, E_7, E_8, F_4, G_2$, how this
classification underpins the structure of the Standard Model, and how the roots
and weights organise particle multiplets.

We assume the reader is comfortable with linear algebra and knows a little quantum
mechanics. No prior exposure to abstract Lie theory is required.

\section{Basic Concepts: Solvability, Semisimplicity, and the Radical}

\begin{definition}
A \textbf{Lie algebra} over $\C$ is a vector space $\g$ with a bilinear bracket
$[\cdot,\cdot]:\g\times\g\to\g$ that is alternating ($[X,Y]=-[Y,X]$) and satisfies
the Jacobi identity
\[
[X,[Y,Z]] + [Y,[Z,X]] + [Z,[X,Y]] = 0.
\]
\end{definition}
Physicists encounter Lie algebras as the commutator algebras of observables. For
instance, the angular momentum operators satisfy $[J_i,J_j] = i\hbar \epsilon_{ijk} J_k$,
which is a real Lie algebra (after factoring out $i$). We will work over the
complex numbers for mathematical convenience; real forms can then be studied by
restricting to suitable real subalgebras.

The \textbf{derived series} is $\g^{(0)} = \g$, $\g^{(k+1)} = [\g^{(k)},\g^{(k)}]$, and the
\textbf{lower central series} is $\g^{[0]} = \g$, $\g^{[k+1]} = [\g,\g^{[k]}]$.
$\g$ is \textbf{solvable} if $\g^{(n)} = 0$ for some $n$, and \textbf{nilpotent}
if $\g^{[n]} = 0$.
These notions formalise the idea of a Lie algebra that can be ``integrated''
step by step. In classical mechanics, solvable Lie algebras appear as algebras of
canonical transformations for which the equations of motion can be solved by
quadratures. In quantum mechanics, a set of operators whose commutators become
successively simpler (nilpotent) generates a Heisenberg–Weyl type algebra.

\begin{example}[Differential equations, motivation]
A finite‑dimensional Lie algebra of vector fields on a manifold $M$ with
$[V_i,V_j] = \sum c_{ij}^k V_k$ is completely integrable (in the sense of Frobenius)
iff it is a subalgebra.  If the Lie algebra is solvable, the corresponding linear
ordinary differential equations can be integrated by successive quadratures —
Lie’s Theorem below guarantees a common eigenvector which reduces the order.
This historical observation ties the algebraic notion of solvability to the
solvability of differential equations by elementary functions.
\end{example}

\begin{definition}
The \textbf{radical} $\rad(\g)$ is the unique maximal solvable ideal of $\g$.
$\g$ is \textbf{semisimple} if $\rad(\g) = 0$, and \textbf{simple} if it has no
proper nonzero ideals and $\dim\g > 1$.
\end{definition}
In physics, semisimple algebras are those for which all solvable ``tails'' have
been stripped away; they are the building blocks of symmetry groups. The Standard
Model gauge group $SU(3)\times SU(2)\times U(1)$ is reductive (direct sum of a
semisimple part $SU(3)\times SU(2)$ and an abelian $U(1)$). The classification
of semisimple algebras thus directly determines possible gauge groups.

\section{Lie’s and Engel’s Theorems (Toolbox)}

The two fundamental theorems that govern solvable and nilpotent Lie algebras are
stated here for completeness; their proofs can be found in \cite{Hall}.

\begin{theorem}[Lie’s Theorem]
If $\g$ is a solvable Lie algebra and $\rho:\g\to\End(V)$ is a representation on
a finite‑dimensional complex vector space $V\neq0$, then there exists a nonzero
$v\in V$ that is a common eigenvector for all $\rho(X)$.
\end{theorem}

The proof inducts on $\dim\g$: pick a codimension‑one ideal $\mathfrak{a}$,
use induction to get an eigenvector for $\mathfrak{a}$, then study the cyclic
space generated by a complementary element $Y$.  The upshot is
\[
[\g,\g]\ \text{acts by nilpotent operators in any representation.}
\]

Physically, this means that if we have a set of operators that generate a solvable
Lie algebra, we can find a basis in which they are all upper‑triangular. This is
the algebraic reason why certain quantum‑mechanical problems (e.g., the harmonic
oscillator) admit a ladder of states that terminates.

\begin{theorem}[Engel’s Theorem]
A Lie algebra $\g$ is nilpotent iff $\ad X$ is nilpotent for every $X\in\g$.
\end{theorem}

The nontrivial direction follows from the existence of a nonzero central element
when all $\ad X$ are nilpotent, and an induction on dimension.

\section{The Killing Form and Cartan’s Criteria}

\begin{definition}
The \textbf{Killing form} on $\g$ is the symmetric bilinear form
\[
\K(X,Y) = \tr(\ad X \circ \ad Y).
\]
It is \emph{invariant}: $\K([X,Y],Z) = \K(X,[Y,Z])$.
\end{definition}
For matrix Lie algebras (subalgebras of $\mathfrak{gl}(n,\C)$), the Killing form
is just (up to a constant) the trace of the product: $\tr(XY)$. For $\mathfrak{su}(2)$,
$\tr(\sigma_a\sigma_b)=2\delta_{ab}$, which is proportional to the usual Euclidean
inner product. This metric structure is fundamental: in gauge theories, the Killing
form is used to write kinetic terms for gauge fields, $- \frac{1}{4g^2} \K(F_{\mu\nu},F^{\mu\nu})$.

\begin{theorem}[Cartan’s First Criterion]
$\g$ is solvable iff $\K(X,Y)=0$ for all $X\in\g$ and $Y\in[\g,\g]$.
\end{theorem}
\begin{proof}[Sketch]
If $\g$ is solvable, Lie’s theorem applied to the adjoint representation
triangularizes all $\ad X$; for $Y\in[\g,\g]$, $\ad Y$ is strictly
upper‑triangular, so its product with any $\ad X$ has zero trace.
Conversely, if the condition holds, one shows that $\g$ cannot contain a subalgebra
isomorphic to $\mathfrak{sl}_2(\C)$ (whose Killing form is nondegenerate),
forcing solvability.  Details are in \cite{Hall}, Thm.~5.32.
\end{proof}

\begin{theorem}[Cartan’s Second Criterion]
$\g$ is semisimple iff $\K$ is nondegenerate.
\end{theorem}
\begin{proof}
Let $\mathfrak{r} = \{X\in\g\mid \K(X,\g)=0\}$ be the radical of the form.
Invariance implies $\mathfrak{r}$ is an ideal; Cartan’s first criterion shows
$\mathfrak{r}$ is solvable, hence $\mathfrak{r}\subseteq\rad(\g)$.
If $\g$ is semisimple, $\rad(\g)=0$, so $\mathfrak{r}=0$ and $\K$ is nondegenerate.
Conversely, if $\K$ is nondegenerate, any solvable ideal $\mathfrak{a}$ satisfies
$\K(\mathfrak{a},\g)=0$ (Lie’s theorem again), so $\mathfrak{a}=0$; thus
$\rad(\g)=0$.
\end{proof}

\begin{corollary}[Decomposition into simples]
A semisimple Lie algebra $\g$ decomposes as a direct sum $\g = \bigoplus \g_i$
of simple ideals, orthogonal with respect to $\K$.
\end{corollary}
\begin{proof}
For any ideal $\mathfrak{a}$, its orthogonal complement $\mathfrak{a}^\perp$
is also an ideal and $\g = \mathfrak{a}\oplus\mathfrak{a}^\perp$ when $\K|_{\mathfrak{a}}$
is nondegenerate.  Iterating yields the decomposition.
\end{proof}

From now on we assume $\g$ is a **complex semisimple Lie algebra**.
The nondegenerate Killing form will be our main tool.
Physically, this tells us that any semisimple gauge algebra splits into
mutually non‑interacting simple factors, each with its own coupling constant.

\section{Cartan Subalgebras, Roots, and the Inner Product}

\subsection{Cartan subalgebras}
\begin{definition}
A \textbf{Cartan subalgebra} $\h$ of $\g$ is a maximal abelian subalgebra consisting
of $\ad$-diagonalizable (semisimple) elements.
\end{definition}
In physical terms, $\h$ corresponds to a maximal set of commuting observables
that can be simultaneously diagonalized. For $SU(2)$, $\h$ is the one‑dimensional
algebra generated by $J_3$; for $SU(3)$, $\h$ is the two‑dimensional algebra
spanned by the two diagonal Gell‑Mann matrices $\lambda_3$ and $\lambda_8$.

Cartan subalgebras always exist: take a \emph{regular} element $H_0\in\g$, i.e.,
one for which the nilspace of $\ad H_0$ has minimal dimension. Then the generalized
nullspace (the span of generalized eigenvectors with eigenvalue $0$) is a Cartan
subalgebra. All Cartan subalgebras are conjugate by inner automorphisms, and their
common dimension is the \textbf{rank} of $\g$. The rank is the number of independent
additive quantum numbers (like isospin $I_3$, hypercharge $Y$) that can be simultaneously
assigned to particles in a representation.

Since $\h$ consists of commuting semisimple operators, they can be simultaneously
diagonalized:
\[
\g = \h \oplus \bigoplus_{\alpha\in R} \g_\alpha,
\tag{4.1}
\]
where $R\subset\h^*\setminus\{0\}$ is a finite set of \textbf{roots}, and
\[
\g_\alpha = \{ X\in\g \mid [H,X] = \alpha(H)X\ \forall H\in\h\}.
\]
The finiteness of $R$ follows because $\g$ is finite‑dimensional. The decomposition
is orthogonal with respect to the Killing form: $\K(\g_\alpha,\g_\beta)=0$ unless
$\alpha+\beta=0$.

The restriction of $\K$ to $\h$ is nondegenerate (otherwise an $H\in\h$ orthogonal
to $\h$ would be orthogonal to all $\g_\alpha$ as well, forcing $H=0$ by nondegeneracy
of $\K$ on $\g$). Thus we have an isomorphism
\[
\phi:\h \longrightarrow \h^*,\qquad H \mapsto \big(H' \mapsto \K(H,H')\big).
\]
For $\lambda\in\h^*$, denote by $t_\lambda$ the unique element of $\h$ with
$\lambda(H) = \K(t_\lambda, H)$ for all $H$.

\subsection{Inner product on the root space}
Define a bilinear form on $\h^*$ by
\[
(\lambda,\mu) := \K(t_\lambda, t_\mu).
\]
Let $E = \Span_\R R \subset \h^*$; this is a real form of $\h^*$ ($\h^* = E\oplus iE$).

\begin{proposition}
$(\cdot,\cdot)$ restricted to $E$ is a real‑valued, positive definite inner product.
\end{proposition}
\begin{proof}
For any $H = t_\lambda$ with $\lambda\in E$, we compute using (4.1) that
\[
\K(H,H) = \tr((\ad H)^2) = \sum_{\alpha\in R} \alpha(H)^2.
\]
Indeed, on $\h$, $\ad H=0$; on each one‑dimensional space $\g_\alpha$, $\ad H$ acts as
multiplication by $\alpha(H)$. Hence the trace of $(\ad H)^2$ is the sum over all roots
of $\alpha(H)^2$. (Note that we will prove $\dim\g_\alpha=1$ shortly, but this formula
holds even before that if we count each root with multiplicity.)
Thus
\[
(\lambda,\lambda) = \K(t_\lambda,t_\lambda) = \sum_{\alpha\in R} \alpha(t_\lambda)^2
= \sum_{\alpha\in R} (\lambda,\alpha)^2 \ge 0.
\]
If $(\lambda,\lambda)=0$, then $(\lambda,\alpha)=0$ for all $\alpha\in R$, so $\lambda$
is orthogonal to the entire span of the roots, i.e., to $E$. Since $(\cdot,\cdot)$ is
nondegenerate on $\h^*$, this forces $\lambda=0$. Hence the form is positive definite.
\end{proof}

Henceforth we write $\langle \lambda,\mu\rangle = (\lambda,\mu)$ for the inner product
on $E$. This Euclidean structure is what allows us to draw root diagrams and measure
angles and lengths. Physically, the inner product gives the metric on the space of
quantum numbers.

\subsection{$\mathfrak{sl}_2$ subalgebras and the fundamental root properties}
Fix a root $\alpha\in R$. Because $\K$ pairs $\g_\alpha$ and $\g_{-\alpha}$ nondegenerately
(if $\alpha$ is a root, so is $-\alpha$, and otherwise $\K(\g_\alpha,\g_\beta)=0$ unless
$\beta=-\alpha$), we can pick $X_\alpha\in\g_\alpha$, $Y_\alpha\in\g_{-\alpha}$ such that
$\K(X_\alpha,Y_\alpha)=1$. A direct calculation shows
\[
[X_\alpha,Y_\alpha] = -t_\alpha.
\]
Indeed, for any $H\in\h$, using invariance,
\[
\K([X_\alpha,Y_\alpha],H) = \K(X_\alpha,[Y_\alpha,H]) = -\K(X_\alpha,\alpha(H)Y_\alpha)
= -\alpha(H)\K(X_\alpha,Y_\alpha) = -\alpha(H).
\]
On the other hand, $\K(t_\alpha,H) = \alpha(H)$ by definition. Since $\K$ on $\h$ is
nondegenerate, we conclude $[X_\alpha,Y_\alpha] = -t_\alpha$.

Now rescale:
\[
E_\alpha = \frac{2}{\langle\alpha,\alpha\rangle}X_\alpha,\quad
F_\alpha = Y_\alpha,\quad
H_\alpha = \frac{2}{\langle\alpha,\alpha\rangle}t_\alpha.
\]
Using $\K(t_\alpha,t_\alpha) = \langle\alpha,\alpha\rangle$, one verifies the
$\mathfrak{sl}_2(\C)$ commutation relations:
\[
[H_\alpha, E_\alpha] = 2E_\alpha,\quad [H_\alpha, F_\alpha] = -2F_\alpha,\quad
[E_\alpha, F_\alpha] = H_\alpha.
\]
Thus for each $\alpha\in R$, we obtain a subalgebra $\mathfrak{s}_\alpha \cong \mathfrak{sl}_2(\C)$.

The adjoint action of this $\mathfrak{s}_\alpha$ on the whole $\g$ makes $\g$ a
(finite‑dimensional) $\mathfrak{sl}_2(\C)$-module. The representation theory of
$\mathfrak{sl}_2(\C)$ is elementary and well‑known to physicists from angular
momentum: in any finite‑dimensional representation,
\begin{itemize}
  \item All eigenvalues of $H_\alpha$ are integers.
  \item The eigenspaces (weight spaces) are finite‑dimensional.
  \item If $v$ is an eigenvector of $H_\alpha$ with eigenvalue $\lambda$ (called a
    weight vector), then $E_\alpha v$ (if nonzero) has eigenvalue $\lambda+2$, and
    $F_\alpha v$ has eigenvalue $\lambda-2$.
  \item The set of eigenvalues (weights) of $H_\alpha$ forms an unbroken string
    $\lambda, \lambda-2, \dots, \lambda-2k$ where $k$ is the length of the string,
    symmetric about zero. In particular, the highest and lowest weights are
    negatives of each other.
\end{itemize}
These facts are proven by standard ladder operator arguments, completely analogous
to those for the harmonic oscillator or angular momentum.

We now apply this to $\g$ itself. Consider the $\alpha$-string through a root $\beta$:
\[
\mathfrak{g}(\beta,\alpha) = \bigoplus_{k\in\Z} \g_{\beta+k\alpha}.
\]
This is an $\mathfrak{sl}_2$-submodule of $\g$, because $E_\alpha\in\g_\alpha$,
$F_\alpha\in\g_{-\alpha}$ and $H_\alpha\in\h$ preserve the root spaces as indicated.

On $\g_{\beta+k\alpha}$, $H_\alpha$ acts as the scalar
\[
\frac{2}{\langle\alpha,\alpha\rangle} (\beta+k\alpha)(t_\alpha)
= \frac{2}{\langle\alpha,\alpha\rangle} \big(\langle\beta,\alpha\rangle + k\langle\alpha,\alpha\rangle\big)
= 2\frac{\langle\beta,\alpha\rangle}{\langle\alpha,\alpha\rangle} + 2k.
\]
Since these numbers must be integers (they are eigenvalues of $H_\alpha$ in a
finite‑dimensional module), we deduce the first fundamental fact:
\[
n_{\beta\alpha} := 2\frac{\langle\beta,\alpha\rangle}{\langle\alpha,\alpha\rangle} \in \Z.
\tag{4.2}
\]
These integers are called \textbf{Cartan integers}.

Moreover, the set of $k$ for which $\g_{\beta+k\alpha}\neq0$ is an unbroken string
from $k=-p$ to $k=q$, where $p,q\ge 0$. The highest weight in this string is
$2\frac{\langle\beta,\alpha\rangle}{\langle\alpha,\alpha\rangle} + 2q$, and the lowest
is $-2\frac{\langle\beta,\alpha\rangle}{\langle\alpha,\alpha\rangle} - 2p$. By symmetry,
\[
2\frac{\langle\beta,\alpha\rangle}{\langle\alpha,\alpha\rangle} + 2q
= -\big(-2\frac{\langle\beta,\alpha\rangle}{\langle\alpha,\alpha\rangle} - 2p\big)
\;\Longrightarrow\; p-q = 2\frac{\langle\beta,\alpha\rangle}{\langle\alpha,\alpha\rangle}.
\tag{4.3}
\]

Now consider the special case $\beta=\alpha$. The string $\g_{\alpha+k\alpha}$ contains
$\g_\alpha$ at $k=0$. Applying $E_\alpha$ to $\g_\alpha$ gives an element of
$\g_{2\alpha}$, but we claim $2\alpha$ is not a root. Indeed, if $2\alpha$ were a root,
the $H_\alpha$ weight on $\g_{2\alpha}$ would be $4$. The $\alpha$-string through $0$
(i.e., taking $\beta=0$, though $0$ is not a root but we can consider the $\h$ part
as $\g_0$) has weights that must be symmetric integers. The action of $H_\alpha$ on $\h$
is zero, so the string through $0$ has weights $\dots, -2,0,2,\dots$. If $\g_{2\alpha}\neq0$,
the weight $4$ would appear, contradicting the fact that in the string through $0$ the
highest weight cannot exceed the maximal eigenvalue, which is bounded by the dimension.
More rigorously, the $\mathfrak{sl}_2$ submodule generated by $\g_\alpha$ (which is a
highest weight vector because $E_\alpha X_\alpha \in \g_{2\alpha}=0$ by the root property)
must have highest weight $2$ (since $H_\alpha$ acts as $2$ on $\g_\alpha$). The dimension
of this irreducible module is $3$, with weights $\{-2,0,2\}$. The weight space of $2$
is $\g_\alpha$ (one‑dimensional), of $0$ is $\h$ (or a subspace thereof), and of $-2$
is $\g_{-\alpha}$. Hence $\dim\g_\alpha=1$, and no other multiples of $\alpha$ ($\pm2\alpha$,
etc.) can appear. This completes the proof that
\begin{enumerate}
  \item $\dim\g_\alpha = 1$,
  \item the only multiples of $\alpha$ in $R$ are $\pm\alpha$.
\end{enumerate}

Finally, using the string condition (4.3) with $p,q$, we compute the reflection
$s_\alpha$ on $\beta$:
\[
s_\alpha(\beta) = \beta - 2\frac{\langle\beta,\alpha\rangle}{\langle\alpha,\alpha\rangle}\alpha
= \beta - (q-p)\alpha.
\]
Because the string is symmetric, $\beta - (q-p)\alpha$ is exactly $\beta - m\alpha$
where $m = q-p$, and one checks that this belongs to the string (it is the element
at position $-p$ after shifting). Hence $s_\alpha(\beta)\in R$. Thus
\begin{enumerate}
  \setcounter{enumi}{2}
  \item $s_\alpha(R) = R$ for all $\alpha\in R$.
\end{enumerate}

We have now established all the fundamental properties of roots using only the
existence of the Cartan subalgebra, the nondegenerate Killing form, and the
representation theory of $\mathfrak{sl}_2$. These properties will be crystallized
into the abstract definition of a root system.

\subsection{The Weyl group}
\begin{definition}
The \textbf{Weyl group} $\W$ is the subgroup of the orthogonal group $O(E)$ generated by
the reflections $\{s_\alpha \mid \alpha\in R\}$.
\end{definition}
Each $s_\alpha$ permutes the finite set $R$, and $R$ spans $E$; hence the map
$\W \to \text{Perm}(R)$ is injective, and $\W$ is a finite group (a subgroup of the
symmetric group on $|R|$ elements). In physical applications, the Weyl group is the
symmetry of the weight diagram; it organizes weights of representations into orbits.

\section{Root Systems and their Geometry}

The properties we have derived are so restrictive and elegant that we now abstract
them, allowing us to classify all possibilities independently of the original Lie
algebra.

\begin{definition}
A \textbf{reduced root system} is a finite set $R$ of nonzero vectors in a Euclidean
space $E$ (with inner product $\langle\cdot,\cdot\rangle$) satisfying:
\begin{enumerate}[label=(R\arabic*)]
  \item $R$ spans $E$.
  \item For each $\alpha\in R$, the only scalar multiples of $\alpha$ that belong to $R$
    are $\pm\alpha$.
  \item For each $\alpha\in R$, the reflection $s_\alpha(\lambda) = \lambda - 2\frac{\langle\lambda,\alpha\rangle}{\langle\alpha,\alpha\rangle}\alpha$
    leaves $R$ invariant.
  \item For any $\alpha,\beta\in R$, $n_{\beta\alpha} = 2\frac{\langle\beta,\alpha\rangle}{\langle\alpha,\alpha\rangle} \in \Z$.
\end{enumerate}
\end{definition}
We have proved that the roots of a semisimple Lie algebra form a reduced root system.
Conversely, every reduced root system arises from some semisimple Lie algebra (Serre’s
theorem, Section 7). Hence classifying semisimple Lie algebras is equivalent to
classifying reduced root systems.

A root system is \textbf{irreducible} if it cannot be expressed as the orthogonal
disjoint union of two nonempty subsets. The Lie algebra is simple iff its root
system is irreducible.

\subsection{Simple roots and bases}
To work with a root system, we choose a direction.

\begin{definition}
A subset $\D \subset R$ is a \textbf{base} (or system of simple roots) if
\begin{itemize}
  \item $\D$ is a basis of $E$,
  \item every $\beta\in R$ can be written as $\beta = \sum_{\alpha\in\D} c_\alpha \alpha$
    with integer coefficients $c_\alpha$ that are either all $\ge 0$ (positive roots)
    or all $\le 0$ (negative roots).
\end{itemize}
\end{definition}

\begin{proposition}
A base always exists.
\end{proposition}
\begin{proof}
Pick a vector $v\in E$ that is \emph{regular}, i.e., $\langle v,\alpha\rangle \neq 0$
for all $\alpha\in R$. (Such vectors exist because the hyperplanes $\alpha^\perp$ do
not cover $E$.) Declare a root $\alpha$ to be \emph{positive} if $\langle v,\alpha\rangle > 0$,
and negative otherwise. An element of $R^+$ (the set of positive roots) is
\emph{ndecomposable} if it cannot be written as a sum of two other positive roots.
One proves that the set $\D$ of all indecomposable positive roots is a base.
The proof uses the finite generation property and integrality. See \cite{Hall}, Thm.~8.16
for details.
\end{proof}

The Weyl group acts simply transitively on the set of all bases.

Given a base $\D = \{\alpha_1,\dots,\alpha_r\}$ ($r$ is the rank), the
\textbf{Cartan matrix} $A = (A_{ij})$ is defined by
\[
A_{ij} = 2\frac{\langle\alpha_j,\alpha_i\rangle}{\langle\alpha_i,\alpha_i\rangle}
= n_{\alpha_j\alpha_i}.
\]
All information about the root system is encoded in this integer matrix. In particular,
\begin{itemize}
  \item $A_{ii} = 2$ for all $i$.
  \item For $i\neq j$, $A_{ij} \in \{0,-1,-2,-3\}$ because $A_{ij}A_{ji} = 4\cos^2\theta_{ij}$
    with $\theta_{ij}$ the angle between $\alpha_i$ and $\alpha_j$. Since the roots are
    not scalar multiples, $A_{ij}A_{ji}$ can only be $0,1,2,3$.
  \item $A_{ij}=0 \iff A_{ji}=0$.
  \item The Cartan matrix is positive definite ($\det A >0$), which follows from the
    positive definiteness of the inner product.
\end{itemize}

\subsection{Dynkin diagrams}
The Cartan matrix is graphically represented by the \textbf{Dynkin diagram}:
\begin{itemize}
  \item Nodes correspond to the simple roots $\alpha_1,\dots,\alpha_r$.
  \item Two nodes $i\neq j$ are connected by $A_{ij}A_{ji}$ edges (0,1,2,3).
  \item If a multiple edge occurs, we draw an arrow pointing from the longer root
    to the shorter root (i.e., from the node with larger $\langle\alpha_i,\alpha_i\rangle$
    to the one with smaller).
\end{itemize}

The root system is irreducible iff the Dynkin diagram is connected. The classification
of irreducible root systems reduces to the classification of connected Dynkin diagrams.

\subsection{Classification theorem}
\begin{theorem}[Cartan–Killing classification]
Every irreducible reduced root system is isomorphic to exactly one of the following
families, with Dynkin diagrams as shown:
\[
\begin{array}{c|l}
\text{Type} & \text{Dynkin diagram} \\ \hline
A_n\ (n\ge1) & \circ\!-\!\circ\!-\!\cdots\!-\!\circ \quad (n\text{ nodes}) \\
B_n\ (n\ge2) & \circ\!-\!\circ\!-\!\cdots\!-\!\circ\!\Rightarrow\!\circ \\
C_n\ (n\ge3) & \circ\!-\!\circ\!-\!\cdots\!-\!\circ\!\Leftarrow\!\circ \\
D_n\ (n\ge4) & \circ\!-\!\circ\!-\!\cdots\!-\!\circ\text{ with an extra node attached to the penultimate} \\
E_6, E_7, E_8 & \text{exceptional branching diagrams} \\
F_4 & \circ\!-\!\circ\!\Rightarrow\!\circ\!-\!\circ \\
G_2 & \circ\!\Rrightarrow\!\circ \quad\text{(triple edge)}
\end{array}
\]
\end{theorem}
\begin{proof}[Proof outline]
The proof is a beautiful exercise in Euclidean geometry.
\begin{enumerate}
  \item Show that the Coxeter graph (ignore arrows) is a tree. If there were a cycle,
    the sum of angles around the cycle would exceed $2\pi$, forcing the Cartan matrix
    to be non‑positive definite, contradiction.
  \item Bound the degree of any node. Since the simple roots are linearly independent,
    a node connected to $k$ other nodes would imply that these $k$ vectors have a
    nontrivial null linear combination, again using positive definiteness. One finds
    that a node can have at most three edges, and certain patterns (like a triple
    edge plus another edge) are forbidden.
  \item These constraints eliminate all graphs except the known families $A_n$,
    $B_n (=C_n)$, $D_n$, $E_6, E_7, E_8$, $F_4$, and $G_2$.
  \item Finally, the arrow directions are determined by the relative root lengths.
    For $B_n$ and $C_n$, which have a double edge, the arrow distinguishes between
    the algebra where the short root is at the end ($B_n$) or the long root is at the
    end ($C_n$). For $F_4$ and $G_2$, the arrows are fixed by consistency.
\end{enumerate}
A full exposition of this combinatorial argument can be found in \cite{Hall}, Ch.~8.
\end{proof}

Consequently, any complex simple Lie algebra belongs to one of the types $A_n$,
$B_n$, $C_n$, $D_n$, $E_6$, $E_7$, $E_8$, $F_4$, $G_2$. A semisimple Lie algebra
is a direct sum of simple such pieces.

\begin{physremark}
In physics, the $A_n$ series corresponds to $\mathfrak{su}(n+1)$; $B_n$ to
$\mathfrak{so}(2n+1)$; $C_n$ to $\mathfrak{sp}(2n)$; $D_n$ to $\mathfrak{so}(2n)$.
The exceptional groups $E_6, E_7, E_8$ appear in Grand Unified Theories (GUTs)
and string theory; $F_4$ and $G_2$ arise in certain supergravity and M‑theory
compactifications.
\end{physremark}

\section{Representations and Weights (A Physicist’s View)}
\label{sec:rep}
For completeness, we briefly sketch how representation theory ties into the above.
A finite‑dimensional representation of $\g$ on a vector space $V$ is a Lie algebra
homomorphism $\pi:\g\to\End(V)$.  When restricted to a Cartan subalgebra $\h$, the
operators $\pi(H)$ for $H\in\h$ are simultaneously diagonalizable (by preservation
of the Killing form, one can define a Cartan involution, etc.).  Thus $V$ decomposes
into \textbf{weight spaces}
\[
V = \bigoplus_{\mu\in\Phi} V_\mu,\qquad V_\mu = \{v\in V \mid \pi(H)v = \mu(H)v\}.
\]
The weights $\mu$ are elements of $\h^*$, and their inner product with the roots
governs which root operators $\pi(X_\alpha)$ map $V_\mu$ to $V_{\mu+\alpha}$.

The \textbf{highest weight} of an irreducible representation is a weight $\lambda$
such that $V_{\lambda+\alpha}=0$ for all positive roots $\alpha$.  The representation
is uniquely determined by this highest weight, which must be \textbf{dominant integral}
(i.e., $2\langle\lambda,\alpha\rangle/\langle\alpha,\alpha\rangle$ is a non‑negative
integer for every simple root $\alpha$).  The Weyl group orbit of the highest weight
generates all weights of the representation, and the multiplicity patterns are given
by the celebrated Weyl character formula.

This machinery is exactly what is used to classify elementary particles in GUTs:
the known particles are assembled into an irreducible representation of the gauge
algebra, and the weights (quantum numbers) are computed using the simple roots of
the Dynkin diagram.

\begin{example}[$SU(3)$ quark model]
For $A_2$ ($\mathfrak{su}(3)$), the rank is 2. The simple roots are $\alpha_1,\alpha_2$
with $\angle=120^\circ$ and equal lengths. The fundamental representation $\mathbf{3}$
has highest weight $\Lambda_1 = \frac{2\alpha_1+\alpha_2}{3}$, while the antifundamental
$\overline{\mathbf{3}}$ has highest weight $\Lambda_2 = \frac{\alpha_1+2\alpha_2}{3}$. The
adjoint representation $\mathbf{8}$ consists of the roots themselves plus the two
zero weights from $\h$. The meson and baryon octets, decuplets, etc., are all explained
by tensor products of these basic representations, governed by the Clebsch–Gordan
coefficients derived from the root and weight geometry.
\end{example}

\section{Reconstruction: Serre’s Theorem}

To close the circle, we state that every abstract root system (or Cartan matrix)
indeed comes from a Lie algebra.

\begin{theorem}[Serre’s presentation]
Let $A = (A_{ij})$ be an $r\times r$ Cartan matrix of finite type (i.e., all
principal minors positive).  Define $\g(A)$ as the Lie algebra generated by
$e_i, f_i, h_i$ ($i=1,\dots,r$) subject to the relations
\[
[h_i,h_j]=0,\quad [e_i,f_j] = \delta_{ij}h_i,\quad
[h_i,e_j] = A_{ij}e_j,\quad [h_i,f_j] = -A_{ij}f_j,
\]
and the \emph{Serre relations}
\[
(\ad e_i)^{1-A_{ij}}(e_j)=0,\quad (\ad f_i)^{1-A_{ij}}(f_j)=0 \quad (i\neq j).
\]
Then $\g(A)$ is a finite‑dimensional complex semisimple Lie algebra whose
Cartan matrix (with respect to $\h = \Span\{h_i\}$) is exactly $A$.
\end{theorem}

\begin{physremark}
In physics, the generators $e_i, f_i, h_i$ correspond to raising, lowering, and
diagonal (Cartan) operators, respectively. The Serre relations are the algebraic
expression of the fact that repeated action of raising/lowering operators eventually
vanishes inside a finite‑dimensional representation — ultimately a consequence of
the finite number of states in a multiplet.
\end{physremark}

Thus the classification of complex simple Lie algebras is perfectly mirrored by
the classification of irreducible Dynkin diagrams, and the whole theory is
completely self‑contained.

\section*{Conclusion}
We have developed the complete classification of semisimple Lie algebras, starting
from first principles and building up to the Dynkin diagrams and Serre’s
reconstruction. The route goes through solvability, the Killing form, Cartan
subalgebras, the $\mathfrak{sl}_2$ subalgebra machinery, root systems, and finally
the combinatorial classification. This framework not only organises all possible
continuous symmetries but also provides the computational tools (roots, weights,
Weyl group) necessary to construct and analyze their representations—the very
vocabulary in which much of modern theoretical physics is written.

\begin{thebibliography}{1}
\bibitem{Hall}
B.~C. Hall, \emph{Lie Groups, Lie Algebras, and Representations: An Elementary Introduction},
2nd ed., GTM 222, Springer, 2015.
\end{thebibliography}

\end{document}